\scanpage{177}
\mspage{51}

grandes devant $\unsure{\pi|\varpi}$ d'éléments de $M^\pi$, ce qui
implique que \deleted{\unsure{tout}} facteur direct d'un
élément de $M^\pi$ est dans $M$, donc aussi
dans $M^\pi$ d'après la condition de stabilité
d) sur $\mathrm{Filt}^{\pi}$. C'est là une condition
restrictive sur la partie $M^\pi$ de $\mathrm{Filt}^{\pi}$.

Aussi, on devra avoir une condition
sur $\unsure{\mathcal M|M}$ plus forte que la validité de
(5.1.) seule :

(5.3.) Pour \inserted{tout $X\in\mathrm{Ob}\,\mathcal M$ et pour} toute partie $J$ de $\mathrm{Ss}_\pi(X)$, telle
que a) $J$ est cofinale dans $\mathrm{Ss}_\pi(X)$ et

% slanted
\margin{Au lieu de $J$ cofinale dans $\mathrm{Ss}_{\pi}(X)$, il suffit $=$ que
l'on ait $\underset{\alpha\in J}{\underrightarrow{\mathrm{lim}}}\,X_\alpha\overset{\sim}{\longrightarrow}X$,
-- la condition corr. sur $M$ devient donc plus exigeante. Mais
si les objets de $\mathrm{Filt}^{\pi}$ sont $\pi$-accessibles, cette condition
est en fait équivalente à la condition plus faible.}
b) Si \deleted{\unsure{on a}} $J'\subset J$ est une partie filtrante
de $J$, de cardinal $\leq\pi$, alors sa borne
sup. $\alpha_{J'}$ dans $\mathrm{Ss}_\pi(X)$ appartient à $J$,

\deleted{l'ens. des $J^{M\deleted{\unsure{\pi}}}$ des $\alpha$}
On a $X\in M$ \unsure{sss|ssi} $J^M\overset{\mathrm{déf}}{=}J\cap\mathrm{Ss}^{M}_\pi(X)$
est cofinal dans $J$ (ou ce qui revient
au $=$, dans $\mathrm{Ss}_\pi(X)$).

\deleted{\unsure{Compte tenu}} Si on admet déjà \unsure{établi|établie} le ``il suffit'' de
(5.1.), alors il suffit dans (5.3.) de
supposer $\deleted{\unsure{M}}X\in M$ et exiger qu'alors
$J^M$ soit cofinal dans $J$.

Mais notons le

% archive sheet number: 176
